\section{Recognition: $X$ is diffeomorphic to $S^6$.}

In this section we prove that the smooth $6$-manifold underlying $X$ is diffeomorphic to the standard Euclidean $6$-sphere $S^6$.

\begin{theorem}[Differentiable Recognition of the 6-Sphere]\label{thm:recognition-S6-full}
The compact complex $3$-manifold $X$, regarded as a smooth closed $6$-dimensional manifold, is smoothly diffeomorphic to the standard Euclidean $6$-sphere:
\[
X \cong_{\mathrm{diff}} S^6.
\]
Consequently, transporting the complex structure of $X$ along this diffeomorphism produces an integrable complex structure on the standard $6$-sphere.
\end{theorem}

\begin{proof}
The proof proceeds in four classical steps:
\begin{enumerate}[label=\textbf{Step \arabic*:}, leftmargin=*]
    \item \textbf{Homotopy Sphere Property:}
    By Corollary~\ref{cor:simple-connectivity}, $\pi_1(X) \cong 0$.
    By Theorem~\ref{thm:homology-consensus}, the integral homology matches $S^6$:
    \[
    H_k(X; \dbZ) \cong H_k(S^6; \dbZ) \quad \text{for all } k \in \dbZ.
    \]
    By the Hurewicz theorem, $\pi_k(X) = 0$ for $k < 6$ and $\pi_6(X) \cong \dbZ$.
    By Whitehead's theorem, any degree-$1$ map $\phi \colon X \to S^6$ is a homotopy equivalence.
    Thus $X$ is a homotopy $6$-sphere.

    \item \textbf{Smale's $h$-Cobordism Theorem:}
    By Smale's classification of high-dimensional manifolds (1962), any smooth homotopy $n$-sphere with $n \ge 5$ is homeomorphic to $S^n$, and its diffeomorphism class is classified by the Kervaire--Milnor group of homotopy spheres $\Theta_n$.

    \item \textbf{Kervaire--Milnor Classification in Dimension 6:}
    By the Kervaire--Milnor theorem (1963), the group of homotopy spheres $\Theta_6$ fits into the exact sequence:
    \[
    0 \to bP_7 \to \Theta_6 \to \pi_6^S / \operatorname{im}(J) \to 0.
    \]
    In dimension $6$, $bP_7 = 0$ because $6$ is even.
    Furthermore, the stable $6$-stem is $\pi_6^S \cong \dbZ/2$, generated by the Toda bracket $\nu^2$.
    The image of the $J$-homomorphism $\operatorname{im}(J) \subset \pi_6^S$ is the entire group $\dbZ/2$.
    Therefore:
    \[
    \Theta_6 \cong \pi_6^S / \operatorname{im}(J) \cong (\dbZ/2) / (\dbZ/2) = 0.
    \]
    There are \textbf{no exotic $6$-spheres}: every smooth homotopy $6$-sphere is diffeomorphic to the standard Euclidean $6$-sphere $S^6$.

    \item \textbf{Conclusion:}
    $X$ is diffeomorphic to standard $S^6$.
    Transporting the complex structure tensor $J \in \mathrm{End}(TX)$ along the diffeomorphism $\psi \colon X \xrightarrow{\sim} S^6$ defines an almost-complex structure $J_{S^6} = d\psi \circ J \circ d\psi^{-1}$ on $S^6$.
    Since $J$ is an integrable complex structure on $X$ ($N_J \equiv 0$), $J_{S^6}$ is an integrable complex structure on the standard smooth $6$-sphere $S^6$, resolving Heinz Hopf's 1947 problem.
\end{enumerate}
\end{proof}
